\section{\texorpdfstring{Static CPI for $r=1$}{Static CPI for r=1}}\label{sec:static}

We prove the central invariant of the paper: at $r{=}1$, vertex peeling never deletes an edge among the surviving vertices, so the stored clique $\Vh(\TPath)\cup\Vp(\TPath)$ of every path remains structurally valid throughout the peel, and the residual state that \cbnd{} maintains (Section~\ref{sec:baseline}) is unnecessary.

\stitle{Static structure and dynamic counters.}
Recall from Section~\ref{sec:prelim} that a path stores the hold set $\Vh(\TPath)$, the pivot set $\Vp(\TPath)$, and the pivot requirement $\eta_{\TPath}=s-|\Vh(\TPath)|$; none of these three objects is edited during peeling.
%
The dynamic state is a single integer $p_{\TPath}$, the number of pivots in $\Vp(\TPath)$ that have not yet been peeled.
%
If a peeled vertex is a hold, every encoded $s$-clique contains that vertex, and the path contributes nothing after the event; if a peeled vertex is a pivot, the path remains the same symbolic object and $p_{\TPath}$ decreases by one.
%
The next theorem records this counter form: a peel event either kills a path contribution or reduces an effective pivot count.

\begin{theorem}[Vertex Removal on \cpi, Counter Form]\label{thm:vertex-removal}
Let $\TPath$ be a path of the \cpi{} of $G$, and let $v\in V$ be removed from $G$.
\begin{enumerate}[leftmargin=1.5em,itemsep=1pt,topsep=2pt]
  \item If $v\in\Vh(\TPath)$, $\TPath$ encodes no surviving $s$-clique.
  \item If $v\in\Vp(\TPath)$, the surviving $s$-cliques are encoded by the same $\Vh(\TPath)$ together with $\Vp(\TPath)\setminus\{v\}$.
%
  Their count depends only on $|\Vp(\TPath)|$.
%
  A counter $p_{\TPath}=|\Vp(\TPath)|$ decremented by one is therefore support-equivalent to deleting $v$ from $\Vp(\TPath)$.
%
  If $p_{\TPath}-1<\eta_{\TPath}$, no surviving $s$-clique is encoded.
  \item Otherwise, $\TPath$ is unaffected.
\end{enumerate}
The dynamic state of $\TPath$ during peeling reduces to a single integer counter $p_{\TPath}$ alongside the fixed pivot requirement $\eta_{\TPath}$; the path is alive iff $p_{\TPath}\ge\eta_{\TPath}$, a derived predicate rather than separately stored state.
%
The stored sets $\Vh(\TPath)$ and $\Vp(\TPath)$ are never modified.
\end{theorem}

\begin{proof}
$\TPath$ encodes $\mathcal{C}(\TPath)=\{\Vh(\TPath)\cup P'\mid P'\subseteq\Vp(\TPath),\ |P'|=\eta_{\TPath}\}$.
%
If $v\in\Vh(\TPath)$, every encoded clique contains $v$ and none survives.
%
If $v\in\Vp(\TPath)$, the surviving cliques are $\{\Vh(\TPath)\cup P'\mid P'\subseteq\Vp(\TPath)\setminus\{v\},\ |P'|=\eta_{\TPath}\}$, a family of size $\binom{|\Vp(\TPath)|-1}{\eta_{\TPath}}$ that depends on $\Vp(\TPath)$ only through its cardinality.
%
By Lemma~\ref{lem:supp}, every support contribution of $\TPath$ also depends only on this cardinality, so the counter update $p_{\TPath}\!\leftarrow\!p_{\TPath}\!-\!1$ is support-equivalent to physical deletion; the family is empty iff $p_{\TPath}\!-\!1<\eta_{\TPath}$.
%
If $v\notin\Vh(\TPath)\cup\Vp(\TPath)$, no encoded clique mentions $v$ and the family is unchanged.
%
\qed
\end{proof}

\stitle{Boundary of the invariant.}
Theorem~\ref{thm:vertex-removal} is specific to $r=1$.
%
When $r\ge2$, a peel event removes an edge or a higher-order clique from the residual object.
%
That event can break the clique represented by $\Vh(\TPath)\cup\Vp(\TPath)$.
%
In that setting, a path may need to be split, shrunk, or rebuilt.
%
The mutable \cpi{} machinery of~\cite{NuclearCD} remains the appropriate general tool, and this paper does not subsume it.

\stitle{A shared path-local counting kernel.}
Both the initial support of Lemma~\ref{lem:supp} and the per-event support loss below depend on a single path-local quantity: the number of $s$-cliques a path contributes to one of its incident vertices at a given pivot scope.
%
We name this kernel and reuse it throughout.

\begin{definition}[Path-local clique count]\label{def:path-contrib}
For a \cpi{} path $\TPath$, a vertex $u\in\Vh(\TPath)\cup\Vp(\TPath)$, and an integer $q\ge0$, the \emph{path-local clique count} of $u$ at pivot scope $q$ on $\TPath$ is
\[
  \pathContrib(\TPath,u,q)=
  \begin{cases}
    \binom{q}{\eta_{\TPath}}, & u\in\Vh(\TPath), \\
    \binom{q}{\eta_{\TPath}-1}, & u\in\Vp(\TPath),
  \end{cases}
\]
with the convention $\binom{m}{j}=0$ when $j<0$ or $j>m$.
%
The parameter $q$ is the number of pivots on $\TPath$ in scope for $u$.
%
When $u\in\Vh(\TPath)$, $q$ is the alive-pivot count of $\TPath$.
%
When $u\in\Vp(\TPath)$, $q$ is the alive-pivot count of $\TPath$ excluding $u$ itself: $u$ fills one of the $\eta_{\TPath}$ pivot slots, and the remaining $\eta_{\TPath}-1$ are chosen from $q$.
\end{definition}

Lemma~\ref{lem:supp} is then a sum of $\pathContrib$ values: with $p=|\Vp(\TPath)|$, every vertex $v$ has $\sdeg(v)=\sum_{\TPath\ni v}\pathContrib(\TPath,v,q_v(\TPath))$, where $q_v(\TPath)=p$ for $v\in\Vh(\TPath)$ and $q_v(\TPath)=p-1$ for $v\in\Vp(\TPath)$.
%
The next lemma uses the same kernel in differential form: a single path event reduces $p$, and the support lost by every other alive vertex is the corresponding $\pathContrib$ differential.

\stitle{Support loss from one path event.}
Theorem~\ref{thm:vertex-removal} leaves only the pivot cardinality changed; the support lost by every other alive vertex on the path is therefore a $\pathContrib$ differential in this cardinality.

\begin{lemma}[Support Loss for a Path Event]\label{lem:support-delta}
Let $\TPath$ be a live path with pre-event pivot count $p_{\mathrm{old}}\ge\eta_{\TPath}$, and let one peel event either (i) decrease its pivot count to $p_{\mathrm{new}}=p_{\mathrm{old}}-d$ with $d\ge1$ and $p_{\mathrm{new}}\ge\eta_{\TPath}$ (path stays alive), or (ii) kill $\TPath$ entirely (hold peel, or pivot peel with $p_{\mathrm{new}}<\eta_{\TPath}$).
%
For every other alive vertex $u$ incident to $\TPath$, the support lost by $u$ from this event is
\[
  \Delta_{\TPath}(u)=\pathContrib(\TPath,u,q_{\mathrm{old}})-\pathContrib(\TPath,u,q_{\mathrm{new}}),
\]
with $q_{\mathrm{old}}=p_{\mathrm{old}}$, $q_{\mathrm{new}}=p_{\mathrm{new}}$ for $u\in\Vh(\TPath)$, and $q_{\mathrm{old}}=p_{\mathrm{old}}-1$, $q_{\mathrm{new}}=p_{\mathrm{new}}-1$ for $u\in\Vp(\TPath)$.
%
Under (ii) the second term vanishes by the binomial convention, so $\Delta_{\TPath}(u)=\pathContrib(\TPath,u,q_{\mathrm{old}})$.
\end{lemma}

\begin{proof}
By Lemma~\ref{lem:supp}, the contribution of $\TPath$ to $u$'s $s$-clique support is $\pathContrib(\TPath,u,q)$ at any current pivot scope $q$.
%
The event changes only the pivot cardinality from $p_{\mathrm{old}}$ to $p_{\mathrm{new}}$ (Theorem~\ref{thm:vertex-removal}); subtracting after from before gives the stated differential.
%
\qed
\end{proof}

\stitle{Running example.}
Figure~\ref{fig:overview} shows the running example at $s=3$: two $K_4$ blocks share edge $(v_3,v_4)$, and two pendant triangles attach at $v_2$ and $v_6$.
%
The six \cpi{} paths in Table~\ref{tab:cpi-leaves} encode the ten triangles through stored $\Vh,\Vp$ sets and the binomial count $\binom{|\Vp|}{s-|\Vh|}$.
%
The $(1,3)$-core values fall into two bands: $\kappa_3(v_i)\!=\!3$ for $i\!\in\!\{1,\ldots,6\}$ (the two $K_4$ blocks) and $\kappa_3(v_i)\!=\!1$ for $i\!\in\!\{7,\ldots,10\}$ (the pendant triangles); the partition is shown by the node shading in Figure~\ref{fig:overview}.
%
The example exercises both events of Theorem~\ref{thm:vertex-removal}: a path dies when a hold vertex is removed, and a path's pivot counter shrinks when a pivot vertex is removed.

\begin{table}[t]
\centering
\caption{\cpi{} paths for the running example of Figure~\ref{fig:overview} at $s=3$.}
\label{tab:cpi-leaves}
\tiny
\begin{tabular}{cccrr}
\toprule
Path & $\Vh$ & $\Vp$ & $\eta=s-|\Vh|$ & $\binom{|\Vp|}{\eta}$ \\
\midrule
$\TPath_1$ & $\{v_7\}$   & $\{v_6,v_8\}$        & $2$ & $1$ \\
$\TPath_2$ & $\{v_9\}$   & $\{v_2,v_{10}\}$     & $2$ & $1$ \\
$\TPath_3$ & $\{v_1\}$   & $\{v_2,v_3,v_4\}$    & $2$ & $3$ \\
$\TPath_4$ & $\{v_2\}$   & $\{v_3,v_4\}$        & $2$ & $1$ \\
$\TPath_5$ & $\{v_3\}$   & $\{v_4,v_5,v_6\}$    & $2$ & $3$ \\
$\TPath_6$ & $\{v_4\}$   & $\{v_5,v_6\}$        & $2$ & $1$ \\
\midrule
\textbf{Total} & & & & $\mathbf{10}$ \\
\bottomrule
\end{tabular}
\end{table}

\begin{example}\label{ex:peel}
Consider the path $\TPath_3=(\{v_1\},\{v_2,v_3,v_4\})$ in Table~\ref{tab:cpi-leaves}, with $p_{\mathrm{old}}=3$ and $\eta=2$. Peeling pivot $v_2$ triggers Lemma~\ref{lem:support-delta} with $p_{\mathrm{new}}=2$: hold $v_1$ loses $\pathContrib(\TPath_3,v_1,3)-\pathContrib(\TPath_3,v_1,2)=\binom{3}{2}-\binom{2}{2}=2$, and pivots $v_3, v_4$ each lose $\pathContrib(\TPath_3,v_i,2)-\pathContrib(\TPath_3,v_i,1)=\binom{2}{1}-\binom{1}{1}=1$. The stored set $\Vp(\TPath_3)$ is unchanged; only the counter $p_{\TPath_3}$ drops from $3$ to $2$.
\end{example}
